% Source: https://artofproblemsolving.com/community/c3h1840969_geo_proof_help
% Stable identity: aops-topic-1840969-q01
We show that every such midpoint $M$ lies on the circle with diameter $OP$.

Let $XY$ be a chord through $P$ with midpoint $M$.
\begin{itemize}
\item If $M\ne O$, the segment $OM$ joins the centre to the midpoint of the chord. Since $OX=OY$, triangle $OXY$ is isosceles and its median $OM$ is also an altitude, so $OM\perp XY$. If moreover $M\ne P$, then $P$ lies on line $XY$, so $\angle OMP=90^\circ$, and by the converse of Thales' theorem $M$ lies on the circle with diameter $OP$. If $M=P$, then $M$ is an endpoint of that diameter, so again it lies on the circle.
\item If $M=O$, then $M$ is the other endpoint of the diameter $OP$, so it lies on the circle as well. (This happens for the chord through $P$ that is a diameter of $\omega$.)
\end{itemize}

\begin{asy}
usepackage("amsmath");
size(6cm);
defaultpen(fontsize(10pt));
real R=3; pair O=(0,0), P=(2.0,0.9);
draw(circle(O,R),linewidth(0.9));
draw(circle((O+P)/2,abs(P-O)/2),red+dashed);
pair chordmid(real t, pen q){
  pair u=dir(t);
  real bb=dot(P,u), cc=dot(P,P)-R^2;
  pair X1=P+(-bb+sqrt(bb^2-cc))*u, X2=P+(-bb-sqrt(bb^2-cc))*u;
  draw(X1--X2,q);
  return (X1+X2)/2;
}
real[] angs={0,40,135,160};
for(real t:angs){ pair M0=chordmid(t,gray(0.5)); dot(M0,red); }
pair M=chordmid(70,linewidth(0.9));
draw(O--M,blue);
real e=0.22; pair v1=e*unit(O-M), v2=e*unit(P-M);
draw(M+v1--M+v1+v2--M+v2,blue);
draw(O--P,dashed);
dot(M,red);dot(O);dot(P);
label("$O$",O,SW);label("$P$",P+0.32*dir(102));label("$M$",M+0.3*dir(318),red);label("$\omega$",R*dir(225),SW);
label("\mbox{}",(min(currentpicture,user=true)+max(currentpicture,user=true))/2,black); // keeps the text after the figure black when LaTeX typesets the labels inline
\end{asy}

So all the midpoints lie on the circle with diameter $OP$ (dashed in the figure). $\blacksquare$

(Conversely, every point $M$ of that circle is such a midpoint: the chord through $P$ perpendicular to $OM$ has midpoint $M$ (for $M=O$ take the diameter through $P$), so the locus is the whole circle with diameter $OP$. If $P=O$, all the chords through $P$ are diameters and all the midpoints coincide with $O$.)
